\section{Discussion}
\label{sec:disc_battery}

This Section reflects on the contribution on modeling fleet mobility and \ac{e-b} battery status developed in Part~\ref{pt:predictive_operations}. It builds on the modeling setup and empirical analysis presented in Chapter~\ref{ch:mobility_and_battery_modeling_bicing}, and discusses the main implications of the proposed framework for battery-aware operations in docked \acp{bss}. Methodological details and results are provided in Sections~\ref{sec:battery_methodology} and~\ref{sec:battery_results}, respectively.

A central finding of this contribution is the presence of pronounced spatiotemporal variation in station-level battery levels. Battery levels were generally lower in the evenings, reflecting the impact of daily usage, and higher in the mornings after overnight resting periods. Spatially, stations in northern neighborhoods exhibited higher average battery levels, likely attributed to lower usage rates and longer resting times. Additionally, notable correlations were observed between stations' battery levels, their inter-event times, and their distance from other stations.

Building on these observations, the developed distance index provides valuable insights into one of the primary operational challenges faced by electric \acp{bss}: battery level management. Stations with lower distance index values are more prone to accelerated battery depletion, as they tend to receive more frequent trips from nearby origins, leaving less time for recharging between trips. From an operational standpoint, this metric offers \ac{bss} operators a practical tool to identify such stations and prioritize interventions, including more frequent battery replacements, recharging, or strategic bike redistribution, thereby enhancing the system's efficiency and reliability.

Methodologically, the presented bike battery forecasting approach combines a traditional probabilistic framework with heterogeneous data sources, resulting in a multifaceted tool capable of improving the real-time management of micro-mobility services. By coupling a validated battery dynamics function with a Markov-chain mobility model, the framework enables joint predictions about future \ac{e-b} availability and battery status. This is particularly relevant in increasingly electrified fleets, where predicting availability alone can be insufficient: a docked bike with a depleted battery is effectively unavailable. The main value of the approach is therefore its ability to move toward microscopic forecasts that capture not only where bikes are likely to be, but also whether they will be usable upon arrival.

At the same time, the scope and accuracy of the results are conditioned by several data limitations. The available battery data were limited in both resolution and coverage: one dataset had an 8-hour temporal resolution, which complicates robust model training and evaluation, while the 30-minute resolution dataset covered only two weeks and did not include bikes in transit. These constraints reduce the ability to learn fine-grained dynamics and to fully validate predictions under a wide range of operating conditions.

In addition, the lack of detailed trip information restricts the ability to model battery consumption more realistically. Because the exact routes taken by users were unavailable, battery depletion could not be estimated from realized paths, and round trips could not be reconstructed. Further sources of heterogeneity that affect consumption are also unobserved, including bicycle condition (e.g., tire inflation) and rider characteristics (e.g., weight or height). These omissions are common in operational datasets, but they imply that the battery consumption component must rely on simplifying assumptions and calibrated parameters rather than on fully observed on-route behavior.

These limitations also point to clear directions for improving the approach. Access to \ac{gps}-based route information would likely improve power-consumption estimates by allowing the battery function to be applied along realized trajectories, while higher-resolution and longer time series could enhance the battery predictions. Additionally, experimental data collected through ad-hoc studies could improve the model's performance by fine-tuning the parameters, ensuring a more accurate representation of real-world conditions.

Overall, this contribution shows that battery-aware forecasting can support operations by anticipating not only where bikes will be, but also whether they will be usable. Despite the limited resolution and missing route-level information, the results demonstrate that meaningful decision support is feasible with already available \ac{bss} data and would further improve with richer trip observations.
